%% =============================================================================
%% SKELETON - Section 4: Benchmark Factor Models
%% =============================================================================
%% Tables/figures INLINE (working draft for supervisor).
%% Same structure as paper1_sec4_benchmarks.tex.
%% Maps to: 02a, 02b, 02c, 03_subperiod_rolling_analysis
%% =============================================================================


\subsection{Benchmark Specifications}
\label{sec:bench_specs}

We estimate three benchmark spanning regressions from the literature,
each designed to capture different sources of systematic risk in fixed
income strategies. \citet{duarte2007risk} specify 10 factors spanning
equity, bond market, and Treasury returns.
\citet{fung2002risk} use 7 asset-based style factors including
trend-following lookback straddles on bonds, currencies, and commodities.
\citet{brooks2020active} decompose active fixed income performance into
traditional risk premia (term, credit, emerging market debt, volatility,
inflation-linked). All models are estimated via OLS with Newey--West HAC
standard errors (6 lags).

Each benchmark is estimated in two panels. The first uses the original
USD-denominated factors, replicating each framework exactly as specified
by its authors. The second uses euro equivalents as a robustness check,
since all three strategies operate in European markets.
For the Fama--French factors---used directly in
\citet{duarte2007risk} and indirectly in \citet{fung2002risk}---we source
the European portfolios from the Kenneth French data library; because
these are reported in USD, we convert them to EUR following
\citet{gluck2021currency}. For long factors (market excess return):
\begin{equation}\label{eq:gluck_long}
  F_{t}^{\textsc{eur}}
  = \frac{1 + F_{t}^{\textsc{usd}} + r_{f,t}^{\textsc{usd}}}
         {1 + r_{t}^{\textsc{usd/eur}}}
  - 1 - r_{f,t}^{\textsc{eur}},
\end{equation}
and for long-short factors (SMB, HML, UMD):
\begin{equation}\label{eq:gluck_ls}
  \mathit{LS}_{t}^{\textsc{eur}}
  = \frac{\mathit{LS}_{t}^{\textsc{usd}}}
         {1 + r_{t}^{\textsc{usd/eur}}},
\end{equation}
where $F_{t}^{\textsc{usd}}$ and $\mathit{LS}_{t}^{\textsc{usd}}$ are
the original factors in US dollars, $r_{f,t}^{\textsc{usd}}$ and
$r_{f,t}^{\textsc{eur}}$ are the US and euro risk-free rates
(Fama--French T-bill rate and Euribor 1M, respectively), and
$r_{t}^{\textsc{usd/eur}}$ is the USD/EUR exchange rate return.
All remaining factors use native European equivalents: European bank
stocks, euro-area corporate bond indices, and German government bonds
(Duarte); euro-area term, credit, and inflation-linked indices
(Brooks et al.); 10Y Bund yield and Baa--Bund credit spread
(Fung \& Hsieh). The Fung \& Hsieh trend-following lookback straddle
factors (PTFSBD, PTFSFX, PTFSCOM) are currency-independent by
construction (D.~Hsieh, private communication) and are therefore
identical in both panels. The complete factor list is reported in
Appendix~\ref{app:factors}; VIF diagnostics are in Appendix~\ref{app:vif}. 


\subsection{Results and Discussion}
\label{sec:bench_results}

\paragraph{Factor Model Regressions.}
Tables~\ref{tab:duarte}--\ref{tab:active_fi} report full
regression results for the three benchmark specifications. Alpha is
statistically significant across all frameworks and strategies, while
$R^2_{\mathrm{adj}}$ remains low (5--15\%), indicating that standard risk
factors do not span returns.

\input{\tablepath/Duarte_article_monthly}

\input{\tablepath/FungHsieh_article_monthly}

\input{\tablepath/ActiveFI_article_monthly}

\paragraph{Cross-Model Comparison and Regime Analysis.}
Table~\ref{tab:alpha_synthesis} summarizes alpha across frameworks
(Panel~A) and decomposes it by stress regime (Panel~B). Alpha is
uniformly higher during HIGH stress months (iTraxx Main $> 100$ bps),
consistent with \citet{duffie2010presidential}: when funding constraints bind,
arbitrage capital withdraws, mispricings widen, and alpha increases.
This pattern motivates the cross-strategy analysis in Section~6.

\input{\tablepath/alpha_synthesis_across_models_monthly}

\paragraph{Conditional Alpha.}
Panel~C of Table~\ref{tab:alpha_synthesis} formalises the regime effect
using the \citet{ferson1996measuring} conditional model:
\begin{equation}\label{eq:ferson_schadt}
  r_{i,t} = \alpha_0 + \alpha_1\, z_t
  + \boldsymbol{\beta}' \mathbf{X}_t + \varepsilon_t,
\end{equation}
where $\mathbf{X}_t$ are the benchmark factors and $z_t$ is the
standardized iTraxx Main 5Y spread, proxying for intermediary funding
stress. Unlike the binary regime split in Panel~B, this specification
allows alpha to vary continuously with market conditions.
The coefficient $\alpha_1 > 0$ indicates that alpha increases during
stress, consistent with \citet{duffie2010presidential}: when
intermediary capital is scarce, arbitrage positions offer a higher
premium to compensate for the additional funding risk and
balance-sheet cost of entering the trade.

\paragraph{Rolling Alpha.}

Figure~\ref{fig:rolling_alpha_composite} plots rolling 36-month alpha
estimates for all three strategies under each benchmark specification.
Alpha spikes visibly during stress episodes (shaded red), confirming
that the regime effect documented in Table~\ref{tab:alpha_synthesis}
is not driven by a single sub-period. Detailed sub-period splits
and threshold robustness for each framework are reported in
Appendix~\ref{app:subperiod}.

\begin{figure}[H]
\centering
\includegraphics[width=\textwidth]{\figpath/rolling_alpha_composite_monthly.pdf}
\caption{Rolling 36-Month Alpha Across Benchmark Factor Models.
  Each panel shows one strategy. Lines represent rolling alpha under
  Duarte et al.\ (2007, blue), Brooks et al.\ (2020, orange), and
  Fung \& Hsieh (2004, green). Red shading: HIGH stress regime
  (iTraxx Main 5Y $> 100$ bps).}
\label{fig:rolling_alpha_composite}
\end{figure}